We measure two margins of carbon leakage from the EU Emissions Trading System --- domestic across-firm reallocation in firm-to-firm transactions, and international cross-border reallocation in customs imports --- on the same Belgian micro-data and under the same post-2017 EU ETS price regime. Both margins are quantitatively small in Belgium.

The within-country across-firm test, on the near-universal Belgian B2B panel 2005--2022, returns precisely-zero reduced-form regression coefficients at both the across-supplier (within NACE 4-digit) and across-category margins. Mapped into a CES nest using the empirical within-cell ETS-saturation moment $E[s_{j^*}(1 - s_{E_n})]$ on the Test~H sample, the within-NACE-4d full-sample bound is $\sigma \in [+0.05, +1.17]$ (95\% CI) under our preferred conservative pass-through assumption and $[+0.52, +1.08]$ under full pass-through, with point estimate at the Cobb-Douglas benchmark $\widehat{\sigma} \approx 1.00$ in both cases. The across-category margin gives a wider but consistent bound, $\sigma_{\text{cat}} \in [0.5, 5.0]$ centered at $1.05$. Both margins place $\sigma$ tightly around the Cobb-Douglas benchmark and the full-sample bound rules out, under both pass-through assumptions, elasticities at the magnitudes typically estimated in import-substitution work for differentiated intermediates.

The across-country cross-border test, replicating \citet{CosterMejeanDiGiovanni2024} on the Belgian customs panel 2000--2019, returns Phase~2 and Phase~3 import-share coefficients of $-0.0037$*** and $-0.0024$**, two orders of magnitude smaller than the corresponding French coefficients of $+0.110$*** and $+0.121$*** under the same identification, and pointing in the opposite direction. The result is robust across four control-group choices and six fixed-effect specifications; no specification recovers CMdG-style positive coefficients. Pass-through of carbon costs into regulated-supplier prices is real on the same Belgian panel (peak monthly local-projection coefficient of $+4.08$ at 12 months on the \citet{kanzig2023unequal} structural shock; Phase~2 and Phase~4 differential PPI gaps of $+0.094$ and $+0.113$). The price signal that would induce reallocation is operative; the two reallocation responses both fail to materialise at quantitatively important levels.

Five mechanisms are jointly consistent with the pattern: small typical shock magnitudes (although the right tail cuts against this as the primary explanation), high relational capital in Belgian B2B, regulated-alternative saturation within heavy NACE 4-digit sectors, intra-EU substitution as the natural reaction to a within-EU policy, and the chemicals-cluster non-ETS-saturation channel specific to the Antwerp port. We cannot separately identify these mechanisms with the current data; the joint pattern is consistent with all five operating in the same direction.

\paragraph{Implications.} The free-allocation rationale and the CBAM rationale rest on different empirical premises about the magnitude of leakage at different margins. Our two reduced-form numbers speak directly to those premises in the Belgian setting: the within-country across-firm elasticity is at Cobb-Douglas (so the leakage rationale for free allocation to immobile installations is weaker in Belgium than the policy currently presumes); the cross-border share coefficient is small and wrong-signed (so the cross-border rationale for CBAM is also weaker on the Belgian customs evidence than the literature presumes). Three margins we do not measure --- production relocation, services trade, and post-2019 imports --- could in principle restore the leakage rationale for either policy, and our results do not bound them. The Belgium-vs-France divergence on the same EU ETS regime is itself an empirical fact that the leakage debate has not engaged with: cross-border leakage is heterogeneous across EU member states by roughly a factor of fifty, and country-level industrial composition appears to be a first-order driver. Calibrating EU-wide carbon-policy models on cross-border-leakage estimates from larger countries (CMdG France) would likely overstate aggregate-emissions effects of unilateral carbon pricing in dense-network small open economies.

We do not interpret our measurements as time-invariant or as predictions about reallocation at higher carbon prices. Whether either margin would activate at $\sim$2.5$\times$ the realized 2022 EUA price --- where right-tail pair-shock magnitudes would scale to match the per-input shocks that identify $\sigma \approx 4$ in trade-microdata work --- is a question we cannot answer with current data, and we leave it open. Within the policy-relevant Phase~III/IV window, our estimates provide a direct micro-foundation for the substitution and import-reallocation parameters that discipline quantitative production-network models of carbon policy: both point to small reallocation responses, and the implications for aggregate carbon-pricing counterfactuals are consequential.
